Suppose in particular that @@M0@@; then the @@M1@@-th universal polynomial of @@M2@@ has no components of weight @@M3@@. Now consider the canonical morphism: @@M4@@; composing with @@M5@@ gives a morphism: @@M6@@, which, by the proof of 1.12, factors through @@M7@@ and defines a morphism @@M8@@; we immediately see that the @@M9@@-th universal polynomial of @@M10@@ is the same as that of @@M11@@, that is, the isobaric component of weight @@M12@@ of the @@M13@@-th universal polynomial of @@M14@@.

The same remarks apply to the case @@M15@@: we can define a morphism of functors @@M16@@ (with @@M17@@, where @@M18@@ is the functor defined on @@M19@@ by @@M20@@) analogously to 1.11; we then repeat the reasoning and conclusions of the preceding paragraph.

\sgaStableTarget{sga6.e05.para.1-16}1.16. Let us finish with an example. Suppose that @@M21@@. We define an endomorphism of the functor @@M22@@, denoted by @@M23@@ and called the Todd character, by putting:
@@M24@@
Conditions i) and ii) of Theorem 1.13 are verified immediately, showing that we have indeed defined an endomorphism of @@M25@@.

\sgaStableTarget{sga6.e05.sec.2}\section*{2. Definition of @@M26@@-rings; examples}

In this section, the base ring of Section 1 will be @@M27@@ (except in 2.9); we work in the category @@M28@@ (1.8) and denote by @@M29@@ the restriction of @@M30@@ to @@M31@@ (1.10).

